\section{Regimes, unresolved physics and the route to a design}\label{sec:validity}

\subsection{Scales that determine which equations may be reduced}
Choose a characteristic geometric length $L>0$, slow velocity $U_s$, slow
variation time $T_s$, acoustic velocity amplitude $V_a$, sound speed $c$, and
reference density, viscosity and surface tension $\rho,\mu,\sigma$. These are
representative positive scales, not assigned apparatus values. Let
$\nu=\mu/\rho$ be kinematic viscosity, $D_T$ thermal diffusivity, and
$\Pi_a$ a characteristic acoustic traction. Where gravity is relevant, $g$ is
its magnitude. Define the following dimensionless ratios and times:
\begin{align}
 \mathrm{Ma}_a&=V_a/c, &
 \mathrm{Re}_s&=\rho U_sL/\mu, &
 \mathcal I_s&=\rho L^2/(\mu T_s),\label{eq:scale-ratios-a}\\
 \mathrm{Oh}&=\mu/\sqrt{\rho\sigma L}, &
 \mathrm{Bo}&=|\Delta\rho|gL^2/\sigma, &
 \mathrm{Ca}_a&=\Pi_aL/\sigma,\label{eq:scale-ratios-b}\\
 \tau_\sigma&=\sqrt{\rho L^3/\sigma}, &
 \tau_\mu&=\mu L/\sigma, &
 \tau_\nu&=\rho L^2/\mu,\label{eq:scale-times}\\
 \delta_\nu&=\sqrt{2\nu/\omega}, &
 \delta_T&=\sqrt{2D_T/\omega}.&&\label{eq:boundary-layer-scales}
\end{align}
Here $\mathcal I_s$ measures unsteady inertia relative to viscosity,
$\mathrm{Oh}$ is the Ohnesorge number, $\mathrm{Bo}$ the Bond number, and
$\mathrm{Ca}_a$ an acoustic traction/capillary scale ratio. The three $\tau$
quantities are capillary--inertial, viscous--capillary and momentum-diffusion
times; $\delta_\nu,\delta_T$ are viscous and thermal penetration depths.
Both fluids may supply important scales, so a favorable ratio in one phase
does not justify eliminating the other.

\begin{longtable}{@{}>{\raggedright\arraybackslash}p{34mm}p{111mm}@{}}
\toprule Reduction & Conditions to examine before using it \\\midrule\endhead
Geometric optics & Optical wavelength small relative to relevant curvature
and wavefront scales; intended polarization, dispersion and diffraction
requirements treated separately.\\
Linear acoustics & $\mathrm{Ma}_a\ll1$, small fractional thermodynamic
perturbations, and acoustic particle displacement small relative to local
geometric scales. The eventual mean interface deformation need not be small.\\
Thin acoustic layers & $\delta_\nu,\delta_T$ small compared with relevant wall
curvature, tangential variation and cavity scales for the particular
boundary-layer approximation \cite{bach2018,jorgensen2021}.\\
Cycle averaging & A window satisfying \cref{eq:averaging-separation}; beats,
modulation and optical exposure must also be resolved or justified as negligible.\\
Settled harmonic field & Acoustic ring-down short relative to envelope and
interface changes, with compatible actuator settling.\\
Creeping mean flow & Both $\mathrm{Re}_s\ll1$ and $\mathcal I_s\ll1$ over the
formation and perturbation times being studied. Small convective inertia
alone does not remove capillary oscillations.\\
Hydrostatic target load & Negligible mean flow, compatible tangential balance,
constant density and surface tension, and the declared volume and wetting law.\\
Isothermal optics & Thermally induced shape, wave-transfer and index changes
small compared with the requested tolerances over the full operating time.\\
\bottomrule
\end{longtable}

As an explicitly limited analytic example, consider two inviscid, infinitely
deep fluids separated by a flat horizontal interface, with the denser fluid
below. Let $k_\perp>0$ be an interfacial spatial wavenumber and
$\omega_\Gamma(k_\perp)$ its gravity--capillary angular frequency. Then
\begin{equation}
 \omega_\Gamma^2(k_\perp)
   =\frac{\Delta\rho\,gk_\perp+\sigma k_\perp^3}{\rho_-+\rho_+}.
 \label{eq:inviscid-capillary-limit}
\end{equation}
This is a scale check for fast interface compliance, not the eigenfrequency
of an unspecified chamber. Finite depth, wetting, curvature and viscosity
change its response; the viscous problem has a nontrivial dispersion relation
\cite{denner2016}. In particular, making the carrier much faster than the
relevant interfacial modes is one ingredient in neglecting first-order
capillary corrections to acoustic scattering. The actual coupled boundary
problem determines whether that neglect is accurate enough optically.

\subsection{Physical effects that may set the attainable optical error}
No single hierarchy is guaranteed for all materials and arrangements. The
following effects enter the same governing problem when their consequences
are comparable with the requested optical tolerance:
\begin{enumerate}
\item \emph{Streaming and absorption}: steady recirculation changes viscous
tractions, temperature and the acoustic field. It is not equivalent to adding
a fitted normal-pressure offset.
\item \emph{Thermal and composition feedback}: density, viscosity, sound speed,
attenuation, surface tension and optical index can all depend on temperature
or composition. Surface-tension gradients create additional tangential stress.
\item \emph{Fast optical modulation}: interface motion and index modulation in
\cref{eq:fast-optical-error} may affect an exposure even when the mean shape
is exactly Cartesian.
\item \emph{Wetting and rheology}: hysteresis, evaporation, contamination,
non-Newtonian response or curing alter the state variables or the admissible
equilibria. Each requires its own measured or justified closure.
\item \emph{Actuator and boundary loading}: radiation, elastic compliance,
transducer coupling and source calibration affect both attainable loads and
their shape derivatives. A finite array is not an arbitrary pressure boundary.
\item \emph{Disturbances and resolution}: vibration, parameter drift and finite
sensor bandwidth act on the actual dynamic system. Tolerances must account for
their optical effect, not only a deterministic design residual.
\end{enumerate}
These are model-selection questions, not reasons to include every effect
indiscriminately. The decision is made by scales and sensitivity against a
stated optical requirement.

Finite actuator authority particularly limits arbitrary fine spatial detail.
For a nearly flat surface with a small sinusoidal height component of
wavenumber $k_\perp$, the capillary part of the required traction scales as
$\sigma k_\perp^2$ times its height amplitude. Thus fine features require large
curvature loads. Whether the available wave fields can provide those loads is
determined by \cref{eq:quadratic-load}, not by the number of pixels in a target
image. Cartesian surfaces form a restricted optical family, which is a useful
design focus, but their membership in that family does not prove actuatability.

\subsection{A theoretical specification before any numerical experiment}
The formulation gives the following ordered construction:
\begin{enumerate}
\item Fix $(n_o,z_o,n_i,z_i)$, optical wavelength, aperture and path level or
vertex. Use the unsquared Cartesian condition and transmission tests to select
the correct surface. Supply any preceding optical system.
\item Complete the fluid boundary outside the clear aperture and check mass,
volume and wetting compatibility. Derive its required load from
\cref{eq:required-traction,eq:cartesian-traction-explicit}, or retain the full
stationary-flow residual if the quiescent reduction fails.
\item Specify a calibrated source variable and physical boundary operators.
Derive the wave responses and quadratic load forms, with a justified coherence,
loss and averaging model. Check load feasibility and actuator limits.
\item Determine whether the candidate can be stably held using the complete
wave--fluid derivative and appropriate inertial, thermal and actuator states.
If feedback is required, specify what is measured and what can be controlled.
\item Formulate a formation trajectory from the supplied initial state, with
optical, energy, rate and terminal constraints. Derive its adjoint and use one
realizable robust policy for the admitted uncertainty set.
\item State explicitly whether the output is a driven liquid diopter, a
temporarily useful transient, or a cured solid. Impose the corresponding
retention and final-optics conditions.
\end{enumerate}
These steps define the problem and its derivatives before solving a specific
instance. They do not produce a universal four-parameter voltage formula:
the optical data specify the target, while the material, support, wave
transfer and control system determine its realizability.

When numerical work resumes, the first checks should be independent analytic
limits, conservation and energy identities, derivatives of the complete
coupled residual, and convergence with the same physical command held fixed.
Wave, geometry, fluid and temporal resolutions should be separated. Stability
requires converged perturbation dynamics and optical growth measures in
addition to force balance. Only after those checks should re-optimization on
refined models be interpreted as a new design. Experimental validation still
requires source and material calibration and measured dynamic and optical
responses. No such numerical or experimental validation is claimed here.
